\section{The Simulated Annealing Bisection Algorithm}
\label{sec:algorithm}

\subsection{Notation and Subgraph Setup}

Let $G_v = (V_v, E_v, w_v)$ denote the subgraph at bisection tree node $v$,
where $V_v$ is the set of census tracts assigned to the subtree rooted at $v$,
$E_v \subseteq E$ is the induced edge set, and $w_v$ inherits weights from the
parent graph.
Let $m = |V_v|$ be the number of tracts.
Let $p_i$ denote the population of tract $i \in V_v$, and let
$P_v = \sum_{i \in V_v} p_i$ be the total subgraph population.
Let $k_v \ge 2$ denote the number of districts to be formed in the subtree
rooted at $v$; \sa{} divides this into a bisection $(k_L, k_R)$ with
$k_L + k_R = k_v$.

A \emph{bisection} at node $v$ is a partition $(S, V_v \setminus S)$ of the
tracts into two non-empty, contiguous groups satisfying the population-balance
constraint:
\[
  \left|\sum_{i \in S} p_i - \frac{k_L}{k_v} P_v \right|
  \;\le\;
  \delta \cdot \frac{P_v}{k_v},
\]
where $\delta = 0.005$ (the $\pm 0.5\%$ federal tolerance).
The \emph{edge cut} of a bisection $(S, V_v \setminus S)$ is
\[
  \EC(S) = \sum_{\substack{\{u,v\} \in E_v \\ u \in S,\; v \notin S}} w(u,v).
\]

\subsection{Node-Seed Derivation}

SA requires a per-node random seed to ensure reproducibility.
The seed must be deterministic given the global base seed and the node's
position in the bisection tree, and must differ across nodes so that
sibling nodes explore different parts of the boundary-flip space.

\begin{definition}[SA Node Seed]
\label{def:sa-seed}
Let $\mathtt{base\_seed}$ be the unsigned 64-bit integer global seed, and let
$\mathtt{node\_path}$ be the string representation of the path from the root
to node $v$ (e.g., \texttt{"0/1/0"} for the left-right-left child of the root).
The \emph{SA node seed} is
\[
  s_v = \mathrm{first8}\bigl(
    \mathrm{SHA\text{-}256}(
      \texttt{"SA\_NODE\_"} \;\|\; \mathtt{node\_path}
      \;\|\; \texttt{"\_"} \;\|\; \mathrm{le8}(\mathtt{base\_seed})
    )
  \bigr),
\]
where $\mathrm{le8}(x)$ encodes $x$ as 8 little-endian bytes and
$\mathrm{first8}(h)$ reads the first 8 bytes of hash $h$ as a little-endian
unsigned 64-bit integer.
\end{definition}

The domain-separation prefix \texttt{"SA\_NODE\_"} ensures that the SA seed
differs from any other SHA-256 application of the base seed in the platform
(e.g., the DIA seed formula uses \texttt{"DIA\_SEED\_V1"}).
The node path string makes sibling seeds independent even for the same base
seed.

\subsection{Cooling Schedule Parameters}

\begin{definition}[Initial Temperature]
\label{def:t0}
Let $\EC_0 = \EC(\Pi_{\mathrm{init}})$ be the edge cut of the \metis{}
initialisation.
The \emph{initial temperature} is
\[
  T_0 = \max\!\bigl(1.0,\; \mathtt{t0\_factor} \times \EC_0\bigr),
\]
where $\mathtt{t0\_factor} \in (0,1)$ is a hyperparameter with default $0.01$.
The maximum with $1.0$ handles the degenerate case $\EC_0 = 0$
(a perfectly connected subgraph that requires no cut) and the test-fixture
case $\mathtt{t0\_factor} = 0$ (forces $T_0 = 1.0$ for regression testing).
\end{definition}

\begin{definition}[Step Count and Final Temperature]
\label{def:steps}
The \emph{step count} is $n_{\mathrm{steps}} = \mathtt{steps\_per\_tract} \times m$,
where $m = |V_v|$.
The \emph{final temperature} is $T_{\mathrm{final}} = 10^{-4}$ (fixed).
\end{definition}

The geometric temperature at step $t \in \{0, 1, \ldots, n_{\mathrm{steps}}-1\}$ is
\[
  T(t) = T_0 \times \left(\frac{T_{\mathrm{final}}}{T_0}\right)^{t / n_{\mathrm{steps}}}.
\]
At $t = 0$, $T = T_0$; at $t = n_{\mathrm{steps}}$, $T$ would reach $T_{\mathrm{final}}$.
The ratio $T_0 / T_{\mathrm{final}}$ spans approximately $10^2$ to $10^5$ for
typical subgraphs, giving a rich exploration-to-exploitation transition.

\subsection{Acceptance Probability Analysis}

\begin{proposition}[Initial Acceptance Probability]
\label{prop:acceptance}
Under the default parameters ($\mathtt{t0\_factor} = 0.01$,
$T_{\mathrm{final}} = 10^{-4}$), the probability of accepting a
$+1$ edge-cut move at the \emph{first} SA step is
\[
  P(\text{accept}\; \Delta\EC = +1,\; t=0)
  = \exp\!\left(\frac{-1}{T_0}\right)
  = \exp\!\left(\frac{-1}{\max(1,\; 0.01 \times \EC_0)}\right).
\]
For representative subgraph sizes:
\begin{enumerate}
\item $\EC_0 = 100$ (small leaf node): $T_0 = 1.0$, acceptance $= e^{-1} \approx 37\%$.
\item $\EC_0 = 500$ (mid-tree node): $T_0 = 5.0$, acceptance $= e^{-0.2} \approx 82\%$.
\item $\EC_0 = 1{,}000$ (root-level node): $T_0 = 10.0$, acceptance $= e^{-0.1} \approx 90\%$.
\end{enumerate}
In all cases, the initial temperature permits meaningful exploration beyond
the \metis{} local optimum.
\end{proposition}

\begin{proof}
Substituting $\Delta\EC = 1$ and $T = T_0$ into the Metropolis acceptance
criterion $\exp(-\Delta\EC / T)$ gives the formula directly.
\end{proof}

\begin{remark}
The scaling $T_0 \propto \EC_0$ is the key design choice.
A fixed $T_0$ independent of $\EC_0$ would under-explore large-EC nodes
(high acceptance needed but $T_0$ too small) and over-explore small-EC
leaf nodes (excessive acceptance of bad moves).
Proportional scaling ensures that the initial acceptance probability for a
unit edge-cut move is a universal constant $e^{-1/(\mathtt{t0\_factor} \times \EC_0)} \cdot e^0$,
independent of subgraph size, as long as $\EC_0 \ge 1/\mathtt{t0\_factor}$.
For small subgraphs where $\EC_0 < 100$ (with the default $\mathtt{t0\_factor} = 0.01$),
the floor $\max(1, \cdot)$ in Definition~\ref{def:t0} activates, ensuring
$T_0 \ge 1$ even at nearly-trivial leaf nodes.
\end{remark}

\subsection{Formal Pseudocode}

Algorithm~\ref{alg:sa} gives the complete \sa{} bisection procedure.

\begin{algorithm}[t]
\caption{\sa{}-\textsc{Bisect}($G_v$, $k_L$, $k_R$, $\mathtt{base\_seed}$,
  $\mathtt{node\_path}$, $\mathtt{steps\_per\_tract}$, $\mathtt{t0\_factor}$)}
\label{alg:sa}
\begin{algorithmic}[1]
\State $m \gets |V_v|$
\State $s_v \gets \mathrm{SA\text{-}Node\text{-}Seed}(\mathtt{base\_seed},
       \mathtt{node\_path})$ \Comment{Definition~\ref{def:sa-seed}}
\State $\Pi_{\mathrm{init}} \gets \metis\text{-}\mathrm{Bisect}(G_v, k_L, k_R, s_v)$
       \Comment{warm start}
\State $\EC_0 \gets \EC(\Pi_{\mathrm{init}})$
\State $T_0 \gets \max(1.0,\; \mathtt{t0\_factor} \times \EC_0)$
\State $n_{\mathrm{steps}} \gets \mathtt{steps\_per\_tract} \times m$
\State $T_{\mathrm{final}} \gets 10^{-4}$
\State $\mathrm{plan} \gets \Pi_{\mathrm{init}}$
\State $\mathrm{best} \gets (\EC_0,\; \Pi_{\mathrm{init}})$
\Comment{(edge-cut, plan) pair}
\State $\mathrm{rng} \gets \mathrm{SmallRng::seed\_from\_u64}(s_v)$
\For{$t = 0$ \textbf{to} $n_{\mathrm{steps}} - 1$}
  \State $T \gets T_0 \times (T_{\mathrm{final}} / T_0)^{t / n_{\mathrm{steps}}}$
         \Comment{geometric cooling}
  \State $(\mathtt{tract},\, \mathtt{target}) \gets
         \mathrm{RandomBoundaryTractAndAdjacentDistrict}(\mathrm{plan}, \mathrm{rng})$
  \State $\mathrm{proposed} \gets \mathrm{Flip}(\mathrm{plan},\;
         \mathtt{tract},\; \mathtt{target})$
  \If{$\mathrm{Contiguous}(\mathrm{proposed})$ \textbf{and}
      $\mathrm{Balanced}(\mathrm{proposed})$}
    \State $\Delta\EC \gets \EC(\mathrm{proposed}) - \EC(\mathrm{plan})$
    \If{$\Delta\EC \le 0$ \textbf{or}
        $\mathrm{rng.gen::<f64>()} < \exp(-\Delta\EC / T)$}
      \State $\mathrm{plan} \gets \mathrm{proposed}$
      \If{$\EC(\mathrm{plan}) < \mathrm{best}.0$}
        \State $\mathrm{best} \gets (\EC(\mathrm{plan}),\; \mathrm{plan}.\mathrm{clone}())$
      \EndIf
    \EndIf
  \EndIf
\EndFor
\State $\Pi^* \gets \mathrm{best}.1$
\State \textbf{return} $(\mathrm{Left}(\Pi^*),\; \mathrm{Right}(\Pi^*))$
\Comment{as (left\_set, right\_set)}
\end{algorithmic}
\end{algorithm}

\paragraph{Contiguity and balance checks.}
The contiguity check at line~13 verifies that both pieces remain connected
after the proposed flip.
This is implemented via a BFS; after the flipped tract is removed from its
source district, the BFS begins from the minimum-index tract remaining in
that district (using the sorted tract ordering), which is a deterministic
choice given the tract numbering.
The balance check verifies that the population of each piece satisfies the
$\pm\delta$ tolerance relative to the target allocation $(k_L / k_v) P_v$.
Proposals failing either check are simply rejected (not counted as a step).
In practice, the fraction of rejected proposals is small ($<5\%$) at
boundary tracts, because interior tracts are never proposed (only boundary
tracts can improve the edge cut) and boundary flips to the larger piece
are more likely to maintain contiguity.

\paragraph{Best-ever tracking.}
Lines~19--21 maintain the best-ever plan seen during the run.
This is critical: during the high-temperature phase, SA may accept moves
that worsen $\EC$ significantly, exploring the local neighbourhood.
Without best-ever tracking, those excursions would be wasted.
The final plan (after cooling to $T_{\mathrm{final}}$) is a local optimum
under greedy descent but may be worse than the best plan seen mid-run.
The best-ever plan captures both the exploration benefit (finding new
starting points for descent) and the exploitation benefit (greedy descent
from those points as temperature drops).

\subsection{CLI Invocation and Implementation}

\sa{} bisection is exposed via the \texttt{--structure simulated-annealing}
flag in the \texttt{redist} binary.
The full invocation for North Carolina is:

\begin{verbatim}
redist state --state NC --year 2020 \
  --structure simulated-annealing \
  --sa-steps-per-tract 10 \
  --sa-t0-factor 0.01
\end{verbatim}

The two SA parameters (\texttt{--sa-steps-per-tract}, \texttt{--sa-t0-factor})
are passed through the three-layer compositor's Structure configuration.
When combined with \texttt{--search convergence --convergence-threshold 600},
the SeedCompositor runs \cs{} over SA bisections---each seed tries a full
SA refinement, and the best across all seeds is returned.

\subsection{Computational Complexity}

\begin{proposition}[Runtime Complexity]
\label{prop:runtime}
\sa{} bisection at a node with $m$ tracts requires:
\begin{itemize}
\item One \metis{} call: $O(m \log m)$.
\item $n_{\mathrm{steps}} = \mathtt{steps\_per\_tract} \times m$ SA steps,
  each involving a boundary-tract flip proposal, a contiguity BFS
  ($O(m)$ worst-case), and an edge-cut delta computation ($O(\deg(\mathtt{tract}))$).
\item Worst-case total: $O(\mathtt{steps\_per\_tract} \times m^2)$ for
  contiguity checks, which dominates.
\end{itemize}
For census-tract graphs, which are planar, the BFS terminates in
$O(\sqrt{m})$ expected time (by the planar graph BFS bound), giving an
empirical cost of $O(\mathtt{steps\_per\_tract} \times m^{3/2})$ per node
in practice.

For a tree with $k-1$ internal nodes averaging $m/\log k$ tracts per node,
the total SA cost per plan is $O(\mathtt{steps\_per\_tract} \times m^{3/2} / \log k)$
in the planar-graph regime.
\end{proposition}

\begin{remark}
The inner loop constant is small (single integer comparisons), and empirical
runtimes in Section~\ref{sec:empirical} confirm the planar-graph regime:
\sa{} with $\mathtt{steps\_per\_tract} = 10$ is approximately $12\times$
slower than a single \metis{} call per node, consistent with $O(m^{3/2})$
rather than $O(m^2)$ scaling, and remains practical for all 50 states.
\end{remark}
